\section{Literature Review}
\label{sec:literature}

\subsection{Existing Research}

Other researchers have already studied parts of this problem, and their work forms the starting point for this thesis.

First, studies have shown that AI can read and extract information from messy, unstructured documents like scanned forms and receipts. However, the accuracy depends a lot on the type of document and the AI model being used. Jaume et al.\ \cite{jaume2019funsd} created a well-known test collection of scanned administrative forms and showed that checking whether the AI got the exact right answer (yes or no) is the best way to measure accuracy for this kind of task. Park et al.\ \cite{park2019cord} did similar work using a collection of receipts and found the same thing. Extracting data correctly from real-world documents is genuinely difficult, and a proper test is needed to measure it.

Second, researchers have shown that large AI models can be made much smaller without losing much of their ability. Dettmers et al.\ \cite{dettmers2022llmint8} proved this in a detailed study, showing that a large AI can be compressed to use much less memory and still work well on a regular laptop. Frantar et al.\ \cite{frantar2022gptq} went even further, developing a method that compresses AI models even more efficiently. Both of these studies are important because they prove that running a capable AI on a normal laptop is actually possible, not just in theory, but in practice.

Third, Strubell et al.\ \cite{strubell2019energy} studied how much energy and money it costs to run AI systems and showed a clear way to calculate and compare those costs. This thesis will use a similar approach to compare the cost of running a local AI tool against the cost of an online AI service. Wei et al.\ \cite{wei2022emergent} also studied the general abilities and limits of large AI models, which helps explain why some models perform better than others on specific tasks.

\subsection{Research Gap}

Even though all of these areas have been studied, nobody has put them all together in one place. No study has compared local AI tools and paid online AI services on the same task, at the same time, measuring accuracy, speed, and cost together. This thesis does exactly that. It fills a gap that is both important in research and directly useful in practice.

\subsection{Research Questions}

Based on everything studied above, this thesis will answer these three specific questions:

\begin{itemize}
    \item \textbf{Question 1:} Is there a real, meaningful difference in accuracy between a local AI tool and a paid online AI service when extracting information from the same documents?
    \item \textbf{Question 2:} How much faster or slower is the local AI tool compared to the paid online service?
    \item \textbf{Question 3:} At what number of documents per day does it become cheaper to use the free local AI tool instead of paying for the online service?
\end{itemize}
